% !TeX root = ../../Book4.tex
% 纲要标题：## 第16章　双复形与复合函子谱序列
% 纲要标签：b4:ch:15
% 纲要位置：计划纲要.md，第 2861 行。

\chapter{双复形与复合函子谱序列}
\label{b4:ch:15}
\BookChapterNumber{b4:ch:15}

\begin{chapterabstract}
本章以双复形的两种滤过比较复合函子的导出计算，构造 Cartan--Eilenberg 消解，证明超上同调与 Grothendieck 谱序列，并在群扩张中计算低阶项。
\end{chapterabstract}

\begin{learninggoals}
\begin{itemize}
\item 用双复形的两种逐度有限滤过比较两种取上同调次序。
\item 说明逐项独立消解的不足，构造并比较 Cartan--Eilenberg 消解。
\item 从两套超上同调谱序列证明 Grothendieck 定理，逐项核对内射像条件。
\item 从第一象限中的箭头位置推出五项正合列。
\item 将离散群扩张化为不变量函子的复合，计算一个非零传递映射。
\end{itemize}
\end{learninggoals}

设 \(N\triangleleft G\)，而 \(A\) 是一个 \(G\)-模。元素要被整个 \(G\) 固定，当然须先被 \(N\) 固定；由于 \(N\) 正规，\(G/N\) 又能作用在 \(A^N\) 上，故 \(A^G=(A^N)^{G/N}\)。这一等式本身很简单，但取不变量一般不保满射。如果先把 \(N\)-不变量的缺口算出，再取商群的不变量，与直接计算 \(G\)-上同调是否一致？同一个双复形可以先按一个方向、再按另一方向取上同调；谱序列给出比较这两种计算的逐层方法。

\ProseParagraphBreak
本章的谱序列论证使用\chapref{b4:ch:14}的构造和有限滤过收敛，也调用\chapref{b4:ch:07}的消解比较、\chapref{b4:ch:08}的导出函子及\chapref{b4:ch:10}的零调对象。最后一节另需\chapref{b4:ch:11}的群上同调、循环群计算和正规子群不变量。\subsecref{b4:subsec:1502:01}、\subsecref{b4:subsec:1502:02}中的 Cartan--Eilenberg 消解定义、构造与比较也为\secref{b4:sec:1704}的有界消解模型提供所需工具。本章其余部分则以谱序列比较复合函子的两种计算，并逐项核对假设。

\ProseParagraphBreak
可对照 Weibel\cite[Sections 5.6--5.8; Section 6.8]{Weibel1994HomologicalAlgebra}。超上同调采用其 Cohomology Variant 5.7.9 的有下界版本；本章把该构造所需的消解存在、提升和同伦论证写入正文。群扩张部分仅处理离散群，连续系数的约定见\subsecref{b4:subsec:1104:03}；离散连续模的导出识别见\cref{b4:thm:C-continuous-derived}；附录中的连续低次计算仍保留其系数条件。

% !TeX root = ../../Book4.tex
% 纲要标题：### 16.1 双复形的两个滤过
% 纲要标签：b4:sec:1501
% 纲要位置：计划纲要.md，第 3624 行。

\section{双复形的两个滤过}
\label{b4:sec:1501}

双复形有两种先取同调的次序，两者都应计算同一个总复形。滤过的有界条件保证收敛，使逐行或逐列的正合性可以用于比较总同调。

% 纲要内容：
%
% * 横向和纵向先取同调
% * 总复形的收敛条件
% * 比较两种计算方法
%

\subsection{横向与纵向先取同调}
\label{b4:subsec:1501:01}

本节的双复形均采用\cref{b4:def:double-complex}的反交换约定：\(\mathrm{d}_h\mathrm{d}_v+\mathrm{d}_v\mathrm{d}_h=0\)，总微分为 \(\mathrm{d}_h+\mathrm{d}_v\)。两个滤过所改变的是查看各项的次序，而不是重新改变微分符号。

\begin{theorem}\label{b4:thm:double-complex-spectral-sequences}
设 \(D^{a,b}\) 为交换范畴中的双复形，且存在整数 \(a_0,b_0\)，使 \(D^{a,b}=0\) 在 \(a<a_0\) 或 \(b<b_0\) 时成立。总复形 \(T=\operatorname{Tot}(D)\) 有两套强收敛谱序列
\[
 {}^{\mathrm I}E_2^{p,q}=H_h^p(H_v^q(D))\Longrightarrow H^{p+q}(T),
 \qquad
 {}^{\mathrm {II}}E_2^{p,q}=H_v^p(H_h^q(D))\Longrightarrow H^{p+q}(T).
\]
第二套中第一指标是原纵次数。两套谱序列均对双复形映射自然。
\end{theorem}
\begin{proof}
在 \(T^n\) 上分别取
\[
 F^pT^n=\bigoplus_{a\ge p}D^{a,n-a},\qquad
 G^pT^n=\bigoplus_{b\ge p}D^{n-b,b}.
\]
给定总次数，非零项有限，两种滤过逐度有限。两份伴随分次分别为
\[
 \operatorname{gr}_F^pT^{p+q}\cong D^{p,q},\qquad
 \operatorname{gr}_G^pT^{p+q}\cong D^{q,p}.
\]
对 \(F\)，纵微分保持第 \(p\) 列，横微分进入下一层而在该商中为零，故第零微分为 \(\mathrm{d}_v\)，第一页是 \(H_v^q(D^{p,\bullet})\)。反交换关系保证 \(\mathrm{d}_h\) 把纵余圈映为纵余圈、把纵余边映为纵余边，诱导第一页微分；再取同调即为第一条公式。对 \(G\)，保留的是原第 \(p\) 行，第零微分为 \(\mathrm{d}_h\)，第一页微分由 \(\mathrm{d}_v\) 诱导；因此第二条公式中的 \(p\) 是原纵次数，\(q\) 是原横次数。两种微分在交换指标后仍满足反交换关系，不须另添符号。\cref{b4:thm:finite-filtration-convergence}分别给出强收敛和自然性。
\end{proof}

符号 \(H_h^p(H_v^q(D))\) 指先对每列取上同调，再对所得横向复形取上同调；它一般不等于把原双复形只取一次核与商。两套第二页不必相同，稳定页也未必逐项同构。若 \(\mathcal F,\mathcal G\) 分别是两种滤过在总上同调上诱导的滤过，则
\[
 {}^{\mathrm I}E_\infty^{p,q}\cong\operatorname{gr}_{\mathcal F}^pH^{p+q}(T),\qquad
 {}^{\mathrm {II}}E_\infty^{p,q}\cong\operatorname{gr}_{\mathcal G}^pH^{p+q}(T).
\]
它们描述同一对象的不同滤过，收敛目标相同并不使这些相邻商相同。

\subsection{总复形的收敛条件}
\label{b4:subsec:1501:02}

第一象限只是上一定理最方便的情形。若双复形整体平移后位于第一象限，证明原样成立；若横向只有有限列而纵向任意，两种滤过也都逐度有限，仍可使用同一证明。更一般地，只要每条对角线有限且所选滤过在该总次数有起止位置，\cref{b4:thm:finite-filtration-convergence}也适用。

\ProseParagraphBreak
若一条对角线上有无限多项，先要说明总化采用直和还是直积，再检查穷尽、分离和收敛。\cref{b4:exm:sum-product-total-difference}已经给出纵列全部正合而直和总复形有非零上同调的例子。因此“先取纵上同调得到零”不能脱离收敛条件而推出总上同调为零。本章的超上同调和复合函子构造都使用有下界复形及非负次数消解，每条总对角线有限，避免引入额外的无限乘积正合性假设。

\subsection{双复形的两种同调计算}
\label{b4:subsec:1501:03}

\begin{example}\label{b4:exm:double-two-computations}
设 \(k\) 为域，取\secref{b4:sec:1405}的四项双复形。纵向先取上同调，第一页在 \((0,1)\)、\((2,0)\) 各剩一份 \(k\)，第二阶微分是负恒等映射。若横向先取上同调，上排 \(ka\to kc\) 与下排 \(kb\to ke\) 均为同构，所以另一套第一页已全部为零。两者最终都给出零上同调，但达到结论的页数不同。
\end{example}

在较大的计算中，可以先寻找哪一个方向的消解更容易取同调。例如\chapref{b4:ch:08}的 \(\operatorname{Hom}(P,I)\) 双复形，沿投射消解方向取上同调与沿内射消解方向取上同调分别留下不同的单行或单列，因而得到 Ext 的两种计算。那里已经用有限消元证明了比较，现在可以把同一事实写成两套退化谱序列。双复形谱序列可对照 Weibel\cite[Section 5.6]{Weibel1994HomologicalAlgebra}。

% !TeX root = ../../Book4.tex
% 纲要标题：### 16.2 超同调谱序列
% 纲要标签：b4:sec:1502
% 纲要位置：计划纲要.md，第 3630 行。

\section{超同调谱序列}
\label{b4:sec:1502}

对复形逐项取内射消解，还必须使原微分的提升严格满足平方为零。Cartan--Eilenberg 消解同时照顾项、圈、边界和同调，才能支持超上同调计算。前两小节给出消解的定义、存在性与比较证明，也是\secref{b4:sec:1704}建立有界消解模型所需的部分；本节末的超上同调谱序列另以\chapref{b4:ch:14}为基础。

% 纲要内容：
%
% * 复形上的函子
% * 消解双复形
% * 由对象同调通向总同调
%

\subsection{复形上的函子}
\label{b4:subsec:1502:01}

设 \(F:\mathcal A\to\mathcal B\) 是交换范畴间的加性左正合函子，\(\mathcal A\) 有足够内射对象。对单个对象，取内射消解后作用 \(F\) 已定义 \(R^nF\)。对复形 \(C^\bullet\)，仅逐项独立取内射消解还不够：原来的微分也必须提升，并且要保持平方为零。若分别选择 \(C^a\to I^{a,\bullet}\) 及微分的比较提升 \(\widetilde{\mathrm d}^{\,a}:I^{a,\bullet}\to I^{a+1,\bullet}\)，则 \(\widetilde{\mathrm d}^{\,a+1}\widetilde{\mathrm d}^{\,a}\) 与零链映射都提升原来的零映射 \(\mathrm d_C^{a+1}\mathrm d_C^a\)。由\cref{b4:prop:comparison-homotopy-unique}，二者链同伦；但同伦于零不等于严格为零，因而这些提升未必组成双复形。

\ProseParagraphBreak
解决方法是同时消解项、余圈、余边和上同调，使每行的核与像受同一构造控制。以下只处理有下界的上链复形。先在原微分次数记为 \(a\)、消解次数记为 \(b\) 的交换方格中构造，最后按 \(\mathrm{d}_v=(-1)^a\partial_v\)、\(\mathrm{d}_h=\partial_h\) 转成反交换方格。

\begin{definition}\label{b4:def:ce-resolution}
设 \(\mathcal A\) 是有足够内射对象的交换范畴，\(C\) 是其中的有下界上链复形。给定带增广 \(C^a\to I^{a,0}\) 的交换双复形 \((I^{a,b},\partial_h,\partial_v)\)，其中 \(b\ge0\)。若对每个 \(a\)，以下四个增广纵列都是内射消解：
\[
 C^a\to I^{a,\bullet},\quad
 B^a(C)\to B_h^a(I^{\bullet,\bullet}),\quad
 Z^a(C)\to Z_h^a(I^{\bullet,\bullet}),\quad
 H^a(C)\to H_h^a(I^{\bullet,\bullet}).
\]
并且 \(C\) 在某个次数以下为零时，\(I\) 也在该横次数以下为零，则称 \(I\) 为 \(C\) 的\term[ce-xiaojie]{Cartan--Eilenberg 消解}[Cartan--Eilenberg resolution]。这里 \(B_h,Z_h,H_h\) 分别按横微分取余边、余圈和上同调。
\end{definition}

之所以要求这些纵列的各项都内射，是为了使每一横行中的
\[
 0\to B_h^a\to Z_h^a\to H_h^a\to0,
 \qquad 0\to Z_h^a\to I^{a,b}\to B_h^{a+1}\to0
\]
都分裂：第一列的左端 \(B_h^a\) 内射，第二列的左端 \(Z_h^a\) 内射，故相应包含都有收缩。这使用的是左端的内射性；仅有右端内射不能保证分裂。任意加性函子保持分裂短正合列，横向上同调与 \(F\) 才能交换。

\subsection{消解双复形}
\label{b4:subsec:1502:02}

\begin{lemma}\label{b4:lem:ce-injective-rows}
设 \(\mathcal A\) 是有足够内射对象的交换范畴。每个 \(\mathcal A\) 中的有下界上链复形 \(C\) 都能嵌入一个有相同下界的复形 \(J\)，使 \(J^a,B^a(J),Z^a(J),H^a(J)\) 均内射，且 \(C\to J\) 在项、余边及上同调上均为单态射。任意满足这些内射性条件的有下界复形 \(J\) 都有如下延拓性质：若有下界复形映射 \(X\to Y\) 在项、余边及上同调上均为单态射，则每个复形映射 \(X\to J\) 都能延拓为 \(Y\to J\)。
\end{lemma}
\begin{proof}
记 \(q_C^a:C^a\twoheadrightarrow C^a/B^a(C)\) 为商态射，\(\beta_C^a:C^a\twoheadrightarrow B^{a+1}(C)\) 为微分的像分解中的满态射。余圈包含诱导单态射
\(H^a(C)\hookrightarrow C^a/B^a(C)\)。选择单态射
\(u^a:B^a(C)\hookrightarrow U^a\)、\(v^a:H^a(C)\hookrightarrow V^a\)，其中 \(U^a,V^a\) 内射；对象为零时取零。由内射性，分别将它们延拓为
\[
\alpha^a:C^a\longrightarrow U^a,\qquad
\eta^a:C^a/B^a(C)\longrightarrow V^a.
\]
构造 \(J\) 时要分别保存这一度的余边、这一度的上同调，以及将由微分送到下一度的余边。\(U^a\) 承担 \(B^a(C)\) 的内射嵌入，\(V^a\) 承担 \(H^a(C)\) 的内射嵌入；第三份则必须使用下一度的 \(U^{a+1}\)，才能接收 \(\beta_C^a:C^a\twoheadrightarrow B^{a+1}(C)\) 的像。因此令
\[
J^a=U^a\oplus V^a\oplus U^{a+1},\qquad
\mathrm{d}_J^a=
\begin{pmatrix}
0&0&1_{U^{a+1}}\\0&0&0\\0&0&0
\end{pmatrix}.
\]
这三项的分工可写在微分旁边：
\[
\begin{aligned}
J^a&=\underbrace{U^a}_{B^a(C)\text{ 的嵌入}}
\oplus\underbrace{V^a}_{H^a(C)\text{ 的嵌入}}
\oplus\underbrace{U^{a+1}}_{B^{a+1}(C)\text{ 的嵌入}},\\
\mathrm{d}_J^a(x,y,z)&=(z,0,0)
\in U^{a+1}\oplus V^{a+1}\oplus U^{a+2}.
\end{aligned}
\]
微分只把当前的第三块送到下一度的第一块，因为它记录的正是下一度余边；当前余边与上同调代表都应被微分杀掉，所以其余块为零。再次微分时，这个像已位于第一块，因而自动为零。
矩阵表示有限双积之间的态射，其靶为 \(J^{a+1}\)。由此
\(B^a(J)=U^a\)、\(Z^a(J)=U^a\oplus V^a\)、\(H^a(J)=V^a\)，均内射。以双积的三个投影规定态射
\[
j^a=\begin{pmatrix}
\alpha^a\\ \eta^a q_C^a\\ u^{a+1}\beta_C^a
\end{pmatrix}:C^a\longrightarrow J^a.
\]
由于 \(\alpha^{a+1}\mathrm{d}_C^a=u^{a+1}\beta_C^a\)、
\(q_C^{a+1}\mathrm{d}_C^a=0\) 和
\(\beta_C^{a+1}\mathrm{d}_C^a=0\)，直接比较三个投影就有
\(j^{a+1}\mathrm{d}_C^a=\mathrm{d}_J^a j^a\)。

为证明 \(j^a\) 单，取其核 \(\kappa:K\hookrightarrow C^a\)。第三个投影给出 \(u^{a+1}\beta_C^a\kappa=0\)，由 \(u^{a+1}\) 单得到 \(\mathrm{d}_C^a\kappa=0\)，故 \(\kappa\) 经 \(Z^a(C)\) 因子化。将这一因子投影到 \(H^a(C)\)，第二个投影及 \(v^a\) 的单射性使该投影为零，因此 \(\kappa\) 又经 \(B^a(C)\) 因子化。最后第一个投影在 \(B^a(C)\) 上就是单态射 \(u^a\)，所以该因子为零，\(K=0\)。\(j\) 在余边和上同调上诱导的态射分别是 \(u^a,v^a\)，故也单。若 \(C^a=0\) 对 \(a<m\) 成立，则还有 \(B^m(C)=0\)；因而 \(U^m=0\)，连同较低次数所取的零对象，保证 \(J^{m-1}=0\)，并使 \(J\) 有相同下界。

下面证明延拓性质。设 \(i:X\hookrightarrow Y\) 满足题设。在模范畴中，若 \(x\in X^a\) 的像在 \(Y^a\) 中成为余边，则 \(i(\mathrm d_Xx)=\mathrm d_Yi(x)=0\)，项上的单射性使 \(x\) 为余圈。再由 \(H^a(i)([x])=0\) 和上同调单射性，得到 \([x]=0\)，故 \(x\) 原本就是 \(X\) 中的余边。一般交换范畴中，用拉回来表示这份交。令
\(P=X^a\times_{Y^a}B^a(Y)\)，并记其到 \(X^a\) 的投影为 \(p\)。链映射关系给出
\(i^{a+1}\mathrm{d}_X^ap=0\)；\(i^{a+1}\) 单使 \(p\) 经 \(Z^a(X)\) 因子化。该因子在 \(H^a(Y)\) 中的像为零，而 \(H^a(i)\) 单，故它在 \(H^a(X)\) 中也为零，进一步经 \(B^a(X)\) 因子化。反过来，\(i\) 把 \(B^a(X)\) 映入 \(B^a(Y)\)，所以这两个因子给出 \(P\simeq B^a(X)\) 作为 \(X^a\) 的子对象。核与余核的因子化因而给出
\[
\ker\bigl(X^a/B^a(X)\longrightarrow Y^a/B^a(Y)\bigr)
\simeq P/B^a(X)=0.
\]
这证明所需商态射是单态射。

对于任意满足引理内射性条件的 \(J\)，两条短正合列
\[
\begin{gathered}
0\longrightarrow B^a(J)\longrightarrow Z^a(J)
\longrightarrow H^a(J)\longrightarrow0,\\
0\longrightarrow Z^a(J)\longrightarrow J^a
\longrightarrow B^{a+1}(J)\longrightarrow0.
\end{gathered}
\]
因左端内射而分裂。选择分裂就得到
\[
J^a\simeq B^a(J)\oplus H^a(J)\oplus B^{a+1}(J),
\]
且在这些坐标下，微分仍是刚才的标准矩阵：它投影到 \(B^{a+1}(J)\)，再通过下一次的余边包含映入 \(J^{a+1}\)。因此只需对这种标准式证明延拓。

记 \(U^a=B^a(J)\)、\(V^a=H^a(J)\)。复形态射 \(f:X\to J\) 的中间投影杀掉 \(B^a(X)\)，故因子化为 \(X^a/B^a(X)\to V^a\)；其第三个投影由链映射关系等于 \(f_U^{a+1}\mathrm{d}_X^a\)，其中 \(f_U^a:X^a\to U^a\) 是第一个投影。所求延拓也没有第三个分量的自由选择：由 \(\mathrm d_J^a(x,y,z)=(z,0,0)\)，链映射条件强制 \(\pi_3g^a=\pi_1g^{a+1}\mathrm d_Y^a\)。因此利用 \(U^a,V^a\) 的内射性，沿 \(X^a\hookrightarrow Y^a\) 及刚证明的商单态射延拓前两个投影，分别记为 \(g_U^a\) 和 \(g_V^a\)，再定义
\[
g^a=\begin{pmatrix}
g_U^a\\g_V^a q_Y^a\\g_U^{a+1}\mathrm{d}_Y^a
\end{pmatrix}:Y^a\longrightarrow J^a.
\]
第三投影的规定与中间投影杀掉余边的性质保证
\(g^{a+1}\mathrm{d}_Y^a=\mathrm{d}_J^ag^a\)。三个投影在 \(X^a\) 上都恢复 \(f^a\)，故 \(g\) 就是所需延拓。
\end{proof}

\begin{theorem}\label{b4:thm:ce-resolution}
设 \(\mathcal A\) 是有足够内射对象的交换范畴。每个 \(\mathcal A\) 中的有下界上链复形都有 Cartan--Eilenberg 消解。对任意两个这样的复形 \(C,D\) 及其消解，复形映射 \(C\to D\) 可提升为消解之间的交换双复形映射；两种提升诱导的总复形映射，在作用任意到交换范畴的加性函子后仍链同伦。因此总化后的上同调，以及\cref{b4:thm:hypercohomology-spectral-sequences}中两套谱序列从第二页起的自然识别，都不依赖消解选择。
\end{theorem}
\begin{proof}
先考察复形短正合列 \(0\to X\xrightarrow{i}Y\to K\to0\)，假设 \(H^a(i)\) 对每个 \(a\) 都单。由\cref{b4:thm:cohomology-long-exact}，所有连接态射
\(\delta^a:H^a(K)\to H^{a+1}(X)\) 为零，因而有上同调短正合列。将次数 \(a,a+1\) 的项短正合列及其微分组成交换图，蛇引理给出余圈列的连接态射
\(Z^a(K)\to X^{a+1}/B^{a+1}(X)\)。复形关系 \(\mathrm d^2=0\) 使其像落在 \(H^{a+1}(X)\) 这一子对象中；它在 \(B^a(K)\) 上为零，所以下降后正是上同调长正合列的 \(\delta^a\)。于是该态射因子化为
\[
Z^a(K)\twoheadrightarrow H^a(K)
\xrightarrow{\delta^a}H^{a+1}(X)
\hookrightarrow X^{a+1}/B^{a+1}(X),
\]
故为零。这证明 \(Z^a(Y)\to Z^a(K)\) 满，并得到
\[
0\longrightarrow Z^a(X)\longrightarrow Z^a(Y)\longrightarrow Z^a(K)\longrightarrow0.
\]
在模范畴中，这正是说 \(K\) 中的余圈可以提升为 \(Y\) 中的余圈：任取提升的微分落在 \(X\) 中，其上同调类是连接态射的值；该值为零时，减去 \(X\) 中一个适当元素便消去微分。

再对 \(X,Y\) 的两条短正合列
\(0\to Z^a(-)\to(-)^a\to B^{a+1}(-)\to0\)
使用蛇引理。项上的单射使余边映射也单，而刚证明的余圈正合性将两个余核分别识别为 \(Z^a(K)\) 与 \(K^a\)。蛇引理因此给出
\[
0\longrightarrow Z^a(K)\longrightarrow K^a
\longrightarrow\operatorname{coker}\bigl(B^{a+1}(X)\to B^{a+1}(Y)\bigr)
\longrightarrow0.
\]
与 \(B^{a+1}(K)=\operatorname{coker}(Z^a(K)\hookrightarrow K^a)\) 比较，便得到余边短正合列。至此，取余复形同时保持了项、余圈、余边和上同调四类短正合列。

令 \(K_0=C\)。逐次用上一引理选择 \(K_b\hookrightarrow I^{\bullet,b}\)，并令
\(K_{b+1}=\operatorname{coker}(K_b\hookrightarrow I^{\bullet,b})\)。刚才的结论在每一步给出
\[
\begin{aligned}
0&\longrightarrow K_b^a\longrightarrow I^{a,b}
\longrightarrow K_{b+1}^a\longrightarrow0,\\
0&\longrightarrow B^a(K_b)\longrightarrow B_h^a(I^{\bullet,b})
\longrightarrow B^a(K_{b+1})\longrightarrow0,\\
0&\longrightarrow Z^a(K_b)\longrightarrow Z_h^a(I^{\bullet,b})
\longrightarrow Z^a(K_{b+1})\longrightarrow0,\\
0&\longrightarrow H^a(K_b)\longrightarrow H_h^a(I^{\bullet,b})
\longrightarrow H^a(K_{b+1})\longrightarrow0.
\end{aligned}
\]
例如最初两步是 \(0\to C\to I^{\bullet,0}\to K_1\to0\) 与 \(0\to K_1\to I^{\bullet,1}\to K_2\to0\)，纵微分 \(\partial_v^0\) 就是商到 \(K_1\) 后再包含进 \(I^{\bullet,1}\)。一般地，纵微分是复形态射
\(\partial_v^b:I^{\bullet,b}\twoheadrightarrow K_{b+1}\hookrightarrow I^{\bullet,b+1}\)，故与横微分交换。下一次纵微分先商掉 \(K_{b+1}\)，恰好杀掉这次的像，所以相邻复合严格为零。四类短正合列说明各增广纵列均正合，而其中的每项由引理内射，正是定义所需的四类内射消解。若 \(C\) 在 \(a<m\) 为零，上一引理的下界构造使所有 \(I^{\bullet,b},K_b\) 保留该下界。

任意 Cartan--Eilenberg 消解也具有刚才的逐行余核结构。在次数零，定义使增广在四类对象上都单；取余复形后，刚才的分析将其四类对象识别为相应余核。四类纵列的正合性又使这些余核单射到下一行，归纳得到相同的四类短正合列。因此以下比较适用于任意选定的消解。

给定 \(f:C\to D\)，这里用目标横行的内射性，沿源的单态射延拓；投射比较定理则用源的投射性，沿目标的满态射提升。先沿 \(C\hookrightarrow I_C^{\bullet,0}\)，把
\(C\to D\to I_D^{\bullet,0}\) 延拓为一个横向复形态射。它诱导第一余复形之间的态射；再沿 \(K_{C,1}\hookrightarrow I_C^{\bullet,1}\)，将
\(K_{C,1}\to K_{D,1}\to I_D^{\bullet,1}\) 延拓。每一步的嵌入在项、余边和上同调上都单，目标横行满足引理的内射性条件，故可以逐次延拓。余核的泛性质保证相邻延拓与纵微分交换，得到所需交换双复形映射。

为证明同伦唯一性，设两种提升为 \(f^b,g^b:I_C^{\bullet,b}\to I_D^{\bullet,b}\)，并记其差为 \(e^b\)。约定负消解次数的行及微分为零，取 \(s^0=0\)，归纳构造横向复形态射
\(s^b:I_C^{\bullet,b}\to I_D^{\bullet,b-1}\)。次数零的差杀掉增广 \(C\)，故经 \(K_{C,1}\) 因子化；沿 \(K_{C,1}\hookrightarrow I_C^{\bullet,1}\) 延拓后得到 \(s^1\)，使
\(e^0=s^1\partial_{v,C}^0\)。

若已在次数 \(b-1\) 得到
\[
e^{b-1}=\partial_{v,D}^{b-2}s^{b-1}
+s^b\partial_{v,C}^{b-1},
\]
令 \(h^b=e^b-\partial_{v,D}^{b-1}s^b\)。两个提升均与纵微分交换，故
\[
\begin{aligned}
h^b\partial_{v,C}^{b-1}
&=\partial_{v,D}^{b-1}e^{b-1}
-\partial_{v,D}^{b-1}s^b\partial_{v,C}^{b-1}\\
&=\partial_{v,D}^{b-1}\partial_{v,D}^{b-2}s^{b-1}=0.
\end{aligned}
\]
于是 \(h^b\) 杀掉 \(\operatorname{im}\partial_{v,C}^{b-1}=K_{C,b}\)，经余复形 \(K_{C,b+1}\) 因子化。沿
\(K_{C,b+1}\hookrightarrow I_C^{\bullet,b+1}\) 延拓为 \(s^{b+1}\)，便有
\[
e^b=\partial_{v,D}^{b-1}s^b+s^{b+1}\partial_{v,C}^b.
\]
这给出了全部纵同伦等式。

反交换化后，在双次数 \((a,b)\) 定义总同伦
\(S^{a,b}=(-1)^as^{a,b}\)。纵向两项给出
\(\partial_vs+s\partial_v\)，横向两项则为
\((-1)^a(\partial_hs-s\partial_h)=0\)，因为每个 \(s^b\) 都是横向复形态射。因此
\(f-g=\mathrm d_DS+S\mathrm d_C\)。总次数只有有限项，加性函子保持这些等式与有限双积，故作用任意加性函子后仍得到总链同伦。

对恒等复形取两个方向的比较映射，它们的复合与恒等总映射链同伦，因而给出与选择无关的总上同调。同伦 \(s^{a,b}\) 保持横次数 \(a\)，降低纵次数 \(b\)：对列滤过，它保持滤过，纵同伦使两种提升在第一页诱导同一映射；对行滤过，它降低滤过一层，不能直接保证第一页映射相同。不过每个 \(s^b\) 都与横微分交换，因而先取横上同调后，它诱导第一页纵复形的链同伦；再取纵上同调才得到第二页映射相同。以后各页由取同调确定。比较映射的复合性质同样由同伦唯一性保证，故第二页起的识别均自然。
\end{proof}

上面的显式嵌入是在复形层面使用内射性；它避免了逐个消解各项后还要修补 \(\mathrm{d}^2=0\) 的问题。消解本身通常很大，实际计算只会使用其横行分裂和四类纵列的消解性质。

\subsection{对象同调与总同调}
\label{b4:subsec:1502:03}

\begin{definition}\label{b4:def:hypercohomology}
设 \(F:\mathcal A\to\mathcal B\) 为交换范畴间加性左正合函子，\(\mathcal A\) 有足够内射对象，\(C\) 为 \(\mathcal A\) 中的有下界上链复形。取其 Cartan--Eilenberg 消解 \((I^{a,b},\partial_h,\partial_v)\)，以 \(\mathrm d_h=\partial_h\)、\(\mathrm d_v=(-1)^a\partial_v\) 反交换化。对 \(n\in\mathbb Z\)，定义
\[
 \mathbb H^n(F,C)=H^n\operatorname{Tot}(F(I)).
\]
这些对象称为 \(F\) 在 \(C\) 上的\term[chaoshangtongdiao]{超上同调}[hypercohomology]，也记作右超导出函子 \(\mathbb R^nF(C)\)。总化每次只有有限项，且上一比较定理保证选择独立及对复形映射的自然性。
\end{definition}
\symidx{\(\mathbb H^n(F,C),\ \mathbb R^nF(C)\)：函子在复形上的超上同调与右超导出函子}

\begin{theorem}\label{b4:thm:hypercohomology-spectral-sequences}
设 \(F:\mathcal A\to\mathcal B\) 为交换范畴间的加性左正合函子，\(\mathcal A\) 有足够内射对象，\(C\) 为 \(\mathcal A\) 中的有下界上链复形。则有两套自然的强收敛谱序列
\[
 {}^{\mathrm I}E_2^{p,q}=H^p(R^qF(C^\bullet))
 \Longrightarrow\mathbb H^{p+q}(F,C),
\]
\[
 {}^{\mathrm {II}}E_2^{p,q}=R^pF(H^q(C))
 \Longrightarrow\mathbb H^{p+q}(F,C).
\]
若 \(C\) 的每项都 \(F\)-零调，则 \(\mathbb H^n(F,C)\cong H^n(F(C))\) 对每个 \(n\in\mathbb Z\) 成立。对 \(\mathcal A\) 中的有下界上链复形 \(D\)，每个拟同构 \(C\to D\) 都诱导超上同调同构。
\end{theorem}
\begin{proof}
对 \(F(I)\) 使用\cref{b4:thm:double-complex-spectral-sequences}。纵列是 \(C^a\) 的内射消解，故纵向先取上同调给出 \(R^bF(C^a)\)。固定 \(b\) 后，这些对象以 \(R^bF(\mathrm d_C^a)\) 为微分组成横向复形：函子性使相邻复合为零，比较映射保证它正是由 \(\partial_h\) 诱导的微分。再取横上同调得到第一套第二页。纵微分只是通常消解微分乘常数 \((-1)^a\)，其核与像不变，故可直接用同一子商识别上同调。若改用逐项乘 \((-1)^{ab}\) 的复形同构作识别，则横微分相应多出 \((-1)^b\)；沿该横行再次作相同的符号调整便恢复通常微分，所得上同调相同。

横向先取上同调时，每行的余边、余圈和上同调均内射，两个基本短正合列分裂，所以加性 \(F\) 保持它们。微分分解为商到余边再包含进下一项，分裂性使这两张映射作用 \(F\) 后仍分别满、单，因此
\[
 \ker F(\partial_h^a)\cong F(Z_h^a(I^{\bullet,b})),\qquad
 \operatorname{im}F(\partial_h^{a-1})\cong F(B_h^a(I^{\bullet,b})).
\]
这给出典范同构
\[
 H_h^a(F(I^{\bullet,b}))\simeq F(H_h^a(I^{\bullet,b})).
\]
而 \(H^a(C)\to H_h^a(I^{\bullet,\bullet})\) 是内射消解，再取纵上同调便得到 \(R^bF(H^a(C))\)，交换两个指标即为第二式。其纵符号同样通过逐次乘 \((-1)^{ab}\) 识别，不改变结果。下界保证两种滤过逐度有限，故均强收敛。

若各 \(C^a\) 都零调，第一套谱序列只剩 \(q=0\) 一行。把横次数的下界平移至零后，单行退化的边缘同态给出 \(H^n(F(C))\cong\mathbb H^n(F,C)\)。若 \(C\to D\) 拟同构，第二套谱序列在第二页的各位置均为同构。两套消解的横次数各有下界，纵次数非负，因此每个固定总次数仅有有限多个双次数位置；相应滤过逐度有限且两边均强收敛。于是可从这一有限页应用\cref{b4:cor:spectral-comparison-finite}，得到目标的上同调同构。
\end{proof}

特别地，把对象 \(M\) 放在次数零时，\(\mathbb H^n(F,M[0])=R^nF(M)\)；若整个 \(C\) 正合，第二套第二页为零，超上同调全部消失。普通左正合函子却未必保持拟同构：取 \(F=\operatorname{Hom}_{\mathbb Z}(\mathbb Z/2\mathbb Z,-)\)，次数零、一上的 \(C=(\mathbb Z\hookrightarrow\mathbb Q)\) 到 \(D=(\mathbb Q/\mathbb Z)[-1]\) 的商映射是拟同构，但 \(F(C)=0\)，而 \(H^1(F(D))\cong\mathbb Z/2\mathbb Z\)。因此所做的进一步消解是在校正普通逐项函子破坏拟同构的问题，不能由一般左正合性推出上面的横向同调交换。内射上链版本的标准叙述见 Weibel\cite[Cohomology Variant 5.7.9]{Weibel1994HomologicalAlgebra}；相应投射链版本及其存在性、比较和同伦见同书\cite[Definition 5.7.1, Lemma 5.7.2, Exercises 5.7.2--5.7.3]{Weibel1994HomologicalAlgebra}。

% !TeX root = ../../Book4.tex
% 纲要标题：### 16.3 Grothendieck 谱序列
% 纲要标签：b4:sec:1503
% 纲要位置：计划纲要.md，第 3636 行。

\section{Grothendieck 谱序列}
\label{b4:sec:1503}

复合两个左正合函子时，第一个函子的导出群不能一般地丢掉。内射像对第二个函子的零调条件，使这组中间信息组织成收敛到复合导出的谱序列。

% 纲要内容：
%
% * 左正合函子复合
% * 内射映到后一函子零调对象的条件
% * 谱序列及低阶正合列
%

\subsection{左正合函子的复合}
\label{b4:subsec:1503:01}

令 \(F:\mathcal A\to\mathcal B\)、\(G:\mathcal B\to\mathcal C\) 为加性左正合函子，前两个交换范畴有足够内射对象。对 \(M\in\mathcal A\) 取内射消解 \(M\to I^\bullet\)，\(F(I^\bullet)\) 的上同调是 \(R^qF(M)\)。因此它通常不是 \(F(M)\) 的正合消解，而是一个必须保留正次上同调的复形。上一节正好提供了对这种复形再作用 \(G\) 的方法。

\ProseParagraphBreak
第二套超上同调谱序列的第二页已经是 \(R^pG(R^qF(M))\)。尚待解决的只有目标：\(\mathbb H^*(G,F(I^\bullet))\) 是否就是复合函子的导出？这一识别需要各项 \(F(I^a)\) 对 \(G\) 零调，从而第一套谱序列退化。这解释了复合函子定理中看似技术性的对象条件。

\subsection{内射对象的零调像}
\label{b4:subsec:1503:02}

\begin{theorem}[Grothendieck 谱序列]\label{b4:thm:grothendieck-spectral-sequence}
设 \(\mathcal A,\mathcal B,\mathcal C\) 为交换范畴，\(\mathcal A,\mathcal B\) 有足够内射对象。设 \(F:\mathcal A\to\mathcal B\)、\(G:\mathcal B\to\mathcal C\) 为加性左正合函子，并且
\[
 R^pG(F(I))=0\quad(p>0)
\]
对每个内射对象 \(I\in\mathcal A\) 成立。则对每个 \(M\in\mathcal A\)，有自然的第一象限强收敛\term[grothendieck-puxulie]{Grothendieck 谱序列}[Grothendieck spectral sequence]
\[
 E_2^{p,q}=R^pG(R^qF(M))\Longrightarrow R^{p+q}(GF)(M).
\]
\end{theorem}
\begin{proof}
取 \(M\to I^\bullet\) 的内射消解，令 \(C=F(I^\bullet)\)。它位于非负次数，\(H^q(C)=R^qF(M)\)，通常仍有正次上同调，故须用\cref{b4:thm:ce-resolution}再取 \(C\) 的 Cartan--Eilenberg 消解 \(C\to J^{\bullet,\bullet}\)。对 \(G(J)\) 的两种滤过应用\cref{b4:thm:hypercohomology-spectral-sequences}：先取纵上同调用零调像识别总上同调，先取横上同调则保留 \(R^qF(M)\) 的信息。第一套给出
\[
 {}^{\mathrm I}E_2^{a,b}=H^a(R^bG(F(I^\bullet))).
\]
记总上同调为 \(H^n=H^n\operatorname{Tot}(G(J))\)，第一套谱序列的目标滤过为 \(\mathcal F\)。假设使 \(b>0\) 的各项全部为零，而零行是 \(H^a(GF(I^\bullet))\)。底行没有出微分，所有入微分的源也为零，故第二页已经稳定。每个总次数 \(n\) 只有 \((n,0)\) 这一项可能非零；强收敛及目标滤过的有限性因此给出
\[
 \mathcal F^nH^n=H^n,\qquad \mathcal F^{n+1}H^n=0.
\]
于是 \(E_\infty^{n,0}=\mathcal F^nH^n/\mathcal F^{n+1}H^n\) 典范地等于整个目标，底行的边缘同态给出自然同构
\[
 H^n\operatorname{Tot}(G(J))
 \simeq H^n(GF(I^\bullet))=R^n(GF)(M).
\]
第二套则给出
\[
 {}^{\mathrm {II}}E_2^{p,q}=R^pG(H^q(F(I^\bullet)))
 =R^pG(R^qF(M)),
\]
目标是刚识别的同一总上同调。两个消解次数非负，所以滤过逐度有限，这证明强收敛。

对 \(M\to N\)，内射比较定理给出 \(I_M\to I_N\)，再用 Cartan--Eilenberg 比较定理得到双复形映射。第二页是导出函子的自然映射。两种内射提升链同伦，故在 \(H^qF(I)\) 上相同，因而第二页相同；以后各页随之相同。目标的识别来自增广 \(G(C)\to\operatorname{Tot}(G(J))\)，与这些比较映射相容，而 \(GF(I)\) 上的同伦给出通常导出函子的自然态射。于是谱序列及其收敛识别都与选择无关。
\end{proof}

若 \(F\) 是一个正合函子的右伴随，它保持内射对象，自动满足定理条件；这是\cref{b4:prop:right-adjoint-injectives}给出的便于检验的充分条件。保持内射不是必要条件，只须高次 \(R^pG(F(I))\) 消失。在这一内射像条件仍成立的前提下，若 \(F\) 正合，谱序列只剩底行，恢复\cref{b4:thm:derived-composite-exact-first}；若 \(G\) 正合，只剩最左列，恢复\cref{b4:prop:derived-composite-exact-second}。

\subsection{谱序列与低阶正合列}
\label{b4:subsec:1503:03}

低阶五项正合列已在\cref{b4:prop:spectral-five-term}由边缘同态和目标滤过推出。这里将它应用于复合函子的计算。

代入 Grothendieck 谱序列，得到
\[
\begin{aligned}
0&\longrightarrow R^1G(FM)\longrightarrow R^1(GF)(M)
\longrightarrow G(R^1F(M))\\
&\xrightarrow{\mathrm{d}_2}R^2G(FM)\longrightarrow R^2(GF)(M).
\end{aligned}
\]
在模范畴中，给定 \(x\in G(R^1F(M))\)，正合性说明它来自 \(R^1(GF)(M)\) 当且仅当 \(\mathrm d_2(x)=0\)；一般交换范畴中，这一陈述写成前一箭头的像等于 \(\ker\mathrm d_2\)。所以 \(\mathrm d_2\) 记录的是一次导出类能否提升到复合导出类的障碍。最右箭头只到达目标二次对象的深层滤过，不能一般地覆盖其余相邻商。定理及低阶列可对照 Weibel\cite[Cohomological Setup 5.8.1, Grothendieck Spectral Sequence Theorem 5.8.3]{Weibel1994HomologicalAlgebra}；该书函子的字母次序与这里相反，两种表述须按复合次序对应。

% !TeX root = ../../Book4.tex
% 纲要标题：### 16.4 群扩张谱序列
% 纲要标签：b4:sec:1504
% 纲要位置：计划纲要.md，第 3642 行。

\section{群扩张谱序列}
\label{b4:sec:1504}

群扩张的上同调可以先沿正规子群计算，再沿商群计算。商群在中间上同调上的作用与低阶传递映射，都是这一分步计算不可省略的数据。

% 纲要内容：
%
% * Lyndon–Hochschild–Serre 谱序列
% * 不变量函子的复合
% * 五项正合列与具体扩张的解释
% * 选择一个明确的有限循环群扩张写出五项正合列中的各项及传递映射；若使用连续群版本，必须转到附录 C 的连续模假设而非无标记替换
%
% ---
%

\subsection{Lyndon–Hochschild–Serre 谱序列}
\label{b4:subsec:1504:01}

设 \(1\to N\to G\xrightarrow{\pi}Q\to1\) 是离散群的扩张，\(M\) 为左 \(\mathbb ZG\)-模。研究 \(G\) 的上同调，可以先研究正规子群 \(N\)，再研究商群 \(Q\) 对所得上同调的作用。这由 Lyndon--Hochschild--Serre 谱序列\footnote{林登（Roger Conant Lyndon，1917-1988），美国数学家，研究群论与拓扑；塞尔（Jean-Pierre Serre，1926-），法国数学家，研究代数拓扑、代数几何与数论。霍赫希尔德已见本卷前注。}组织。第二步需要保留作用；即使已经算出底交换群 \(H^q(N,M)\)，也还不足以计算商群的上同调。

\begin{theorem}[Lyndon--Hochschild--Serre 谱序列]\label{b4:thm:lhs-spectral-sequence}
设 \(1\to N\to G\to Q\to1\) 是离散群扩张，\(M\) 为左 \(\mathbb ZG\)-模。则有自然的第一象限强收敛\term[lhs-puxulie]{Lyndon--Hochschild--Serre 谱序列}[Lyndon--Hochschild--Serre spectral sequence]
\[
 E_2^{p,q}=H^p(Q,H^q(N,M))\Longrightarrow H^{p+q}(G,M).
\]
其中 \(Q\) 在 \(H^q(N,M)\) 上的作用由 \(G\) 的共轭和系数作用共同诱导。其边缘同态分别为膨胀与限制所诱导的映射。
\end{theorem}

例如在一次上同调中，对余圈 \(c:N\to M\) 和 \(g\in G\)，作用由
\[
 (g\cdot c)(n)=g\cdot c(g^{-1}ng)
\]
给出：输入先共轭，输出再接受系数作用。\cref{b4:prop:normal-subgroup-invariants}已证明 \(N\) 自身的作用在上同调上平凡，故这才下降为 \(Q\) 作用；余链上的作用不必通过 \(Q\)。下一小节给出谱序列的证明。这里所有群都带离散结构；连续群上同调需要连续模及相应消解，不能直接把 \(G\) 换成拓扑群而保留全部假设。

\subsection{不变量函子的复合}
\label{b4:subsec:1504:02}

\begin{proof}[\cref{b4:thm:lhs-spectral-sequence}的证明]
在模范畴之间取
\[
 F:\mathbb ZG\text{-}\mathbf{Mod}\longrightarrow
 \mathbb ZQ\text{-}\mathbf{Mod},\quad F(M)=M^N,
 \qquad G_0=(-)^Q.
\]
\(N\) 正规保证 \(M^N\) 带 \(Q\) 作用，且 \(G_0F(M)=M^G\)。两个函子都加性左正合。

由\cref{b4:prop:normal-subgroup-invariants}，\((- )^N\) 是膨胀函子的右伴随；膨胀不改变底交换群的正合性，因而正合。用\cref{b4:prop:right-adjoint-injectives}，\(F\) 送内射 \(\mathbb ZG\)-模到内射 \(\mathbb ZQ\)-模，特别地其像对 \((- )^Q\) 零调。这一伴随核对的是复合函子定理所需的内射像条件。

取 \(G\)-内射消解 \(M\to I^\bullet\)。还要核对中间导出确实计算 \(N\) 的群上同调，这用的是另一套伴随：\(\mathbb ZG\) 作为右 \(\mathbb ZN\)-模自由，故诱导 \(\mathbb ZG\otimes_{\mathbb ZN}-\) 正合；其右伴随限制函子由同一命题保持内射对象。因此 \(I^\bullet\) 限制到 \(N\) 后仍是内射消解，得到底交换群的自然同构
\[
 R^qF(M)=H^q\bigl((I^\bullet)^N\bigr)\cong H^q(N,M).
\]

还须核对这一同构与 \(Q\) 作用相容。\cref{b4:prop:normal-subgroup-invariants}的证明用双复形
\[
 K^{i,j}=\operatorname{Hom}_{\mathbb ZN}(B_i(N),I^j)\qquad(i,j\ge0)
\]
给出两个增广拟同构
\[
 (I^\bullet)^N\longrightarrow\operatorname{Tot}(K)
 \longleftarrow C^\bullet(N,M).
\]
在 \(K\) 上，\(g\in G\) 对 bar 坐标作共轭，并作用于系数 \(I^j\)。bar 增广在共轭下不变，而 \(M\to I^0\) 是 \(G\)-线性的，所以两个比较映射都与 \(G\) 作用交换。它们诱导的上同调同构因此也是 \(G\) 等变的。左边 \(N\) 逐项平凡作用；右边的共轭余链作用在取上同调后才使 \(N\) 平凡作用，故两边都下降为 \(Q\) 作用。这就把底交换群的识别提升为 \(\mathbb ZQ\)-模的自然同构。

现在 Grothendieck 定理的全部假设已经核对。代入\cref{b4:thm:grothendieck-spectral-sequence}，第二页是所写群上同调，目标为不变量函子 \((- )^G\) 的导出，即 \(H^*(G,M)\)。底行的增广对应把 \(Q\)-余链沿 \(G\to Q\) 拉回并用 \(M^N\hookrightarrow M\) 作为系数映射，这正是膨胀；最左列的增广对应把 \(G\)-余链限制到 \(N\)，并由正规性落在 \(Q\)-不变类中。因此边缘同态是所述膨胀和限制。自然性与强收敛均由复合函子构造得到。
\end{proof}

这个证明没有要求群有限。有限性将在具体消解计算中使用，或者使某些平均算子可用；它不是谱序列存在的必要条件。标准叙述可参阅 Weibel\cite[Lemma 6.8.1, Lyndon/Hochschild-Serre Spectral Sequence 6.8.2]{Weibel1994HomologicalAlgebra}。

\subsection{五项正合列与群扩张}
\label{b4:subsec:1504:03}

由\cref{b4:prop:spectral-five-term}，每个离散群扩张及左 \(\mathbb ZG\)-模 \(M\) 给出
\begin{equation}\label{b4:eq:lhs-five-term}
\begin{aligned}
0&\longrightarrow H^1(Q,M^N)\xrightarrow{\mathrm{inf}}H^1(G,M)
\xrightarrow{\mathrm{res}}H^1(N,M)^Q\\
&\xrightarrow{\mathrm d_2^{0,1}}H^2(Q,M^N)
\xrightarrow{\mathrm{inf}}H^2(G,M).
\end{aligned}
\end{equation}
这里的微分 \(\mathrm d_2^{0,1}\) 称为\term[chuandi-yingshe]{传递映射}[transgression]。这个名称专指谱序列中的低阶微分，不是有限指数子群的 transfer（亦称余限制）。其含义是：一个 \(N\) 上的 \(Q\)-不变一次上同调类，若想成为 \(G\) 上一次类的限制，必须且只须其传递类消失。

\ProseParagraphBreak
若系数 \(A\) 的作用平凡，一次上同调就是群同态 \(G\to A\)。取满足 \(s(1)=1\) 的集合截面 \(s:Q\to G\)，记
\[
 c(x,y)=s(x)s(y)s(xy)^{-1}\in N.
\]
对 \(Q\)-不变同态 \(\phi:N\to A\)，定义 \(u(x,y)=\phi(c(x,y))\)。群乘法结合律给出
\[
 c(x,y)c(xy,z)=s(x)c(y,z)s(x)^{-1}c(x,yz).
\]
施加 \(\phi\)，利用它在共轭下不变，得到
\[
 u(x,y)+u(xy,z)=u(y,z)+u(x,yz),
\]
所以 \(u\) 是归一化 \(2\)-余圈。若 \(f:G\to A\) 延拓 \(\phi\)，令 \(b(x)=f(s(x))\)，则 \(u(x,y)=b(x)+b(y)-b(xy)=\delta b(x,y)\)。反过来，若 \(u=\delta b\)，则在唯一分解 \(g=ns(x)\) 下定义 \(f(g)=\phi(n)+b(x)\)，由共轭不变性及截面乘法公式可验证 \(f\) 为所需同态。因此 \([u]\) 检测延拓障碍。

\ProseParagraphBreak
具体符号还要与谱序列的代表元比较。取本节已经使用的 \(G\) 内射消解 \(A\xrightarrow{j}I^0\to I^1\to\cdots\)，记其微分为 \(\partial\)。\(I^0\) 限制到 \(N\) 后内射，故 \(H^1(N,I^0)=0\)。一次余圈 \(j\phi\) 因而是余边，可选 \(a\in I^0\) 使
\[
na-a=-j\phi(n)\quad(n\in N),\qquad \alpha=\partial a\in(I^1)^N.
\]
这里 \(n\alpha-\alpha=\partial(na-a)=0\)，且 \(\partial\alpha=0\)。在\cref{b4:prop:normal-subgroup-invariants}所用的 \(N\) 的 bar 消解与内射消解的比较双复形中，零次值 \(-a\) 的横、纵微分分别是 \(j\phi\) 与 \(-\alpha\)，所以
\[
\mathrm d(-a)=(j\phi,-\alpha),\qquad
(j\phi,0)-\mathrm d(-a)=(0,\alpha).
\]
因此 \(\alpha\) 对应的是 \([\phi]\)，而不是它的负类。

\ProseParagraphBreak
用 \(Q\) 的 bar 消解计算第二步，双复形可写为
\[
D^{q,p}=C^p\bigl(Q,(I^q)^N\bigr),\qquad
\mathrm d_h=\partial,\qquad \mathrm d_v=(-1)^q\delta_Q.
\]
第一指标 \(q\) 是原消解次数，第二指标 \(p\) 是商群余链次数；谱序列按第二指标滤过，第二页则写作 \(E_2^{p,q}\)。各 \((I^q)^N\) 都是 \(Q\) 内射模，bar 计算与内射计算经消解比较给出前面的复合函子谱序列。以上总化符号与\cref{b4:def:hypercohomology}相同。令 \(\beta_x=s(x)a-a\)。由 \(\phi\) 的共轭不变性，
\[
n\beta_x-\beta_x
=s(x)\bigl((s(x)^{-1}ns(x))a-a\bigr)-(na-a)=0,
\]
所以 \(\beta\in C^1(Q,(I^0)^N)\)，并且 \(\delta_Q\alpha=\partial\beta\)。\(\alpha\) 位于 \((q,p)=(1,0)\)，其商群方向的微分带负号；\(\beta\) 位于 \((0,1)\)，该方向的微分带正号。修正代表元后有
\[
\mathrm d(\alpha+\beta)
=-\delta_Q\alpha+\partial\beta+\delta_Q\beta
=\delta_Q\beta.
\]
由截面的乘法公式，再用 \(A\) 的平凡作用及 \(\phi\) 的共轭不变性，得到
\[
\begin{aligned}
(\delta_Q\beta)(x,y)
&=\bigl(s(x)s(y)-s(xy)\bigr)a\\
&=s(xy)\bigl((s(xy)^{-1}c(x,y)s(xy))a-a\bigr)\\
&=-j\phi(c(x,y))=-j u(x,y).
\end{aligned}
\]
它落在 \(\ker\partial=j(A)\)，正是第二页目标的系数。把每个 \(D^{q,p}\) 分量的元素乘以 \((-1)^{pq}\)，就得到从当前总复形到另一总复形的同构：商群方向的微分 \((-1)^q\delta_Q\) 变为通常的 \(\delta_Q\)，消解方向的微分则变为 \((-1)^p\partial\)。这个同构保持按 \(p\) 的滤过；在第二页的源 \((p,q)=(0,1)\) 和靶 \((2,0)\) 处，它的系数都为一。因此由\cref{b4:thm:filtered-spectral-sequence}的代表元规则，得到
\[
\mathrm d_2^{0,1}[\phi]=-[u].
\]
这里使用了平凡系数，任意系数时不能省去交叉同态的作用项。下一小节在有限循环群中计算这项障碍。低阶正合列的教材叙述见 Weibel\cite[Low Degree Terms 6.8.3]{Weibel1994HomologicalAlgebra}。

\subsection{有限循环群扩张的传递映射}
\label{b4:subsec:1504:04}

取
\[
 1\longrightarrow N=\langle t^2\rangle\longrightarrow
 G=\langle t\mid t^4=1\rangle\longrightarrow Q=C_2\longrightarrow1,
 \qquad M=\mathbb F_2
\]
并令系数作用平凡。所有共轭作用也平凡。由\cref{b4:thm:cyclic-group-cohomology}，五项列中五个非零候选项均为 \(\mathbb F_2\)：
\[
 0\longrightarrow\mathbb F_2\xrightarrow{\mathrm{inf}}
 \mathbb F_2\xrightarrow{\mathrm{res}}\mathbb F_2
 \xrightarrow{\mathrm d_2^{0,1}}\mathbb F_2
 \xrightarrow{\mathrm{inf}}\mathbb F_2.
\]
第一次膨胀把 \(Q\) 的生成元映到一的同态拉回到 \(t\mapsto1\)，所以为恒等同构。限制却为零，因为任意 \(f:G\to\mathbb F_2\) 满足 \(f(t^2)=2f(t)=0\)。正合性已迫使 \(\mathrm d_2^{0,1}\) 为同构，最后一次膨胀为零。

\ProseParagraphBreak
还可以写出传递类的余圈。用加法记 \(Q=\{0,1\}\)，取 \(s(0)=1\)、\(s(1)=t\)。令 \(\phi:N\to\mathbb F_2\) 由 \(\phi(t^2)=1\) 定义，则
\[
 c(x,y)=s(x)s(y)s(x+y)^{-1}\in N,
 \qquad u(x,y)=\phi(c(x,y))
\]
仅在 \((x,y)=(1,1)\) 取一。直接代入余圈等式
\(u(y,z)-u(x+y,z)+u(x,y+z)-u(x,y)=0\)，或者用截面乘法的结合律，即可验证 \(u\) 是\(2\)-余圈。它不是余边：若 \(u=\delta v\)，则 \(0=u(0,y)=\delta v(0,y)=v(0)\)，而 \(\delta v(1,1)=v(1)-v(0)+v(1)=0\)，与 \(u(1,1)=1\) 矛盾。\(H^2(C_2,\mathbb F_2)\) 只有一个非零元，因此已证明非零的传递映射 \(\mathrm d_2^{0,1}\) 把 \([\phi]\) 送到 \([u]\)。本例特征为二，\(-[u]=[u]\)，与上述符号公式一致。

\ProseParagraphBreak
为直接核对最后一次膨胀为零，对 \(0\le i\le3\) 定义一余链 \(b(t^i)=\lfloor i/2\rfloor\pmod2\)。把 \(i=2a+r\)、\(j=2a'+r'\)（\(r,r'\in\{0,1\}\)）代入，得到
\[
 \delta b(t^i,t^j)=b(t^j)-b(t^{i+j})+b(t^i)=rr'
 =u(\pi(t^i),\pi(t^j)).
\]
指数取模四不影响右边的模二计算。因此膨胀后的\(2\)-余圈确实成为余边。与分裂扩张 \(C_2\to C_2\times C_2\to C_2\) 比较，后一情形每个 \(N\to\mathbb F_2\) 都通过第一因子延拓，限制为满，故传递映射为零。这个差别说明低阶微分记录了群扩张本身的信息，而不只由核群与商群的同构类型决定。

\ProseParagraphBreak
要使负号可见，取 \(G=C_9=\langle t\mid t^9=1\rangle\) 和扩张 \(1\to\langle t^3\rangle\to G\to C_3\to1\)，系数为平凡的 \(\mathbb F_3\)。用 \(0,1,2\) 表示商群元素，取 \(s(i)=t^i\)、\(\phi(t^3)=1\)。截面因子集给出进位余圈
\[
u(i,j)=\left\lfloor\frac{i+j}{3}\right\rfloor\pmod3,
\qquad \mathrm d_2^{0,1}[\phi]=-[u].
\]
这两种符号确实代表不同的类。对任意一余链 \(v:C_3\to\mathbb F_3\)，
\[
\sum_{i=0}^{2}(\delta v)(i,1)
=\sum_{i=0}^{2}\bigl(v(1)-v(i+1)+v(i)\bigr)
=3v(1)=0,
\]
其中 \(i+1\) 在 \(C_3\) 中计算。而 \(\sum_i u(i,1)=1\)、\(\sum_i(-u(i,1))=2\)，所以 \([u]\ne-[u]\)。若以 \([\phi]\) 和 \([u]\) 分别作为两端一维空间的基，本书的传递映射就是乘 \(-1=2\) 的同构。


\begin{chaptersummary}
双复形的横纵两个滤过使同一个总上同调接受两种计算。Cartan--Eilenberg 消解把各项及其上同调同时消解，其分裂横行保证加性函子可以与横向上同调交换。由此，一套谱序列先导出逐项函子，另一套先取原复形上同调。

\ProseParagraphBreak
复合函子的内射像条件使第一套谱序列退化并识别目标，第二套才成为 Grothendieck 谱序列。正规子群不变量是正合膨胀函子的右伴随，满足这一条件，因而给出群扩张谱序列。五项正合列中的传递映射检测一次类能否由子群延拓到整个群；四阶循环群的例子说明该映射确实可能非零。
\end{chaptersummary}

\BookExercises

\begin{exercise}\label{b4:ex:15-coprime-double}
把两项上链复形 \(\mathbb Z\xrightarrow{2}\mathbb Z\) 与 \(\mathbb Z\xrightarrow{3}\mathbb Z\) 张量总化，各自位于次数零、一。分别用两个滤过计算其总上同调。
\end{exercise}
\begin{solution}
取横微分为乘二，纵微分在第 \(p\) 列为 \((-1)^p\) 乘三。纵向先取上同调，只在纵次数一留下两项 \(\mathbb Z/3\)，横微分为乘二，是同构，所以第二页全零。横向先取上同调，只在横次数一留下两项 \(\mathbb Z/2\)，纵微分为负乘三，仍是同构，所以第二套第二页也全零。滤过有限，故总上同调全零。符号不同不影响这里的同构性。
\end{solution}

\begin{exercise}\label{b4:ex:15-equal-two-double}
把上题两个乘法都改为乘二，计算总上同调，并直接写出总微分核与像验证。
\end{exercise}
\begin{solution}
先纵向取上同调，只在 \(q=1\) 留下两项 \(\mathbb Z/2\)，横微分为零。因此 \(H^1\cong\mathbb Z/2\)、\(H^2\cong\mathbb Z/2\)、其余为零。直接总化得到
\[
 \mathbb Z\xrightarrow{x\mapsto(2x,2x)}\mathbb Z^2
 \xrightarrow{(a,b)\mapsto-2a+2b}\mathbb Z.
\]
第一映射单，第二映射的核是对角 \(\mathbb Z\)，商掉第一映射的像后为 \(\mathbb Z/2\)；第二映射的余核同样为 \(\mathbb Z/2\)。
\end{solution}

\begin{exercise}\label{b4:ex:15-hyper-shift-object}
设 \(F\) 为本章所讨论的左正合函子，\(M\) 为对象，\(s\in\mathbb Z\)。把 \(M\) 放在次数 \(s\)，即取 \(M[-s]\)。计算其超上同调。
\end{exercise}
\begin{solution}
原复形只在次数 \(s\) 有上同调 \(M\)。第二套超上同调谱序列只在 \(q=s\) 一行非零，\(E_2^{p,s}=R^pF(M)\)，没有可能的高阶微分。每个总次数只有一个相邻商，所以 \(\mathbb H^n(F,M[-s])=R^{n-s}F(M)\)，当 \(n<s\) 时为零。负的 \(s\) 仅平移第一象限，不改变逐度有限性。
\end{solution}

\begin{exercise}\label{b4:ex:15-exact-hyper}
设 \(F\) 正合，\(C\) 为有下界复形。证明 \(\mathbb H^n(F,C)\cong F(H^n(C))\)，并说明不需要选择上同调的代表元。
\end{exercise}
\begin{solution}
正合性使 \(R^pF=0\) 对 \(p>0\) 成立，第二套谱序列只剩 \(p=0\) 列，边缘同态给出自然同构 \(\mathbb H^n(F,C)\cong F(H^n(C))\)。同构由核、像、商及滤过的典范态射给出，没有选择分裂或代表元。也可由各项自动零调先得 \(\mathbb H^n=H^n(F(C))\)，再利用正合函子保持核与像得到同一识别。
\end{solution}

\begin{exercise}\label{b4:ex:15-acyclic-quasi}
设 \(f:C\to D\) 为有下界复形之间的拟同构，且两者每项都对左正合函子 \(F\) 零调。证明 \(F(f)\) 也是拟同构。
\end{exercise}
\begin{solution}
超上同调的第二套谱序列在第二页由 \(R^pF(H^q(f))\) 给出同构，因此其目标为同构。各项零调使自然边缘同态分别识别 \(\mathbb H^n(F,C)=H^n(F(C))\)、\(\mathbb H^n(F,D)=H^n(F(D))\)。由自然性，目标上的同构就是 \(H^n(F(f))\)，故为拟同构。
\end{solution}

\begin{exercise}\label{b4:ex:15-naive-quasi-failure}
令 \(F=\operatorname{Hom}_{\mathbb Z}(\mathbb Z/2,-)\)，\(C=(\mathbb Z\hookrightarrow\mathbb Q)\) 位于次数零、一，\(D=(\mathbb Q/\mathbb Z)[-1]\)。证明自然商映射 \(C\to D\) 拟同构，但 \(F(C)\to F(D)\) 不是拟同构。
\end{exercise}
\begin{solution}
\(C\) 的零次上同调为零，一次为 \(\mathbb Q/\mathbb Z\)，故商映射拟同构。\(F(\mathbb Z)=F(\mathbb Q)=0\)，所以 \(F(C)=0\)。另一方面，在 \(1\bmod 2\) 处取值给出
\[
 F(\mathbb Q/\mathbb Z)\cong
 \{\,x\in\mathbb Q/\mathbb Z\mid 2x=0\,\}
 =\{\,0,\frac12+\mathbb Z\,\}\cong\mathbb Z/2.
\]
因此 \(F(D)\) 在次数一有非零上同调。这里的 \(2\)-挠子群不同于题设中用方括号表示的复形移位。失效条件是 \(\mathbb Z\) 不是 \(F\)-零调对象。经超上同调校正后，两者重新同构，并在次数一得到 \(\mathbb Z/2\)。
\end{solution}

\begin{exercise}\label{b4:ex:15-field-ce}
设 \(k\) 为域，\(C\) 是 \(k\)-向量空间的有下界复形。证明可以把 \(C\) 自身作为仅有消解次数零的 Cartan--Eilenberg 消解，并指出这里与一般模的差别。
\end{exercise}
\begin{solution}
向量空间及其子空间、商空间都内射。因此 \(C^a,B^a(C),Z^a(C),H^a(C)\) 的恒等增广都是长度零的内射消解。把 \(I^{a,0}=C^a\)、其他纵项取零即满足定义。一般环上的模、子模和商模不一定内射，所以把消解次数限制为零通常不满足定义。例如整数群 \(\mathbb Z\) 已不是内射对象。
\end{solution}

\begin{exercise}\label{b4:ex:15-change-rings-ext}
设 \(R\to S\) 为含幺环同态，\(M\) 为左 \(S\)-模，\(A\) 为左 \(R\)-模。证明存在
\[
 E_2^{p,q}=\operatorname{Ext}_S^p(M,\operatorname{Ext}_R^q(S,A))
 \Longrightarrow\operatorname{Ext}_R^{p+q}(M,A).
\]
说明中间对象的左 \(S\) 作用，并核对复合函子的内射像条件。
\end{exercise}
\begin{solution}
取 \(F(A)=\operatorname{Hom}_R(S,A)\)，左 \(S\) 作用为 \((s f)(t)=f(ts)\)。它是正合限制标量函子的右伴随，故保持内射对象。再取 \(G=\operatorname{Hom}_S(M,-)\)，内射对象当然 \(G\)-零调。伴随同构
\[
 \operatorname{Hom}_S(M,\operatorname{Hom}_R(S,A))
 \cong\operatorname{Hom}_R(M,A)
\]
由在 \(1\in S\) 处取值给出，逆映射把 \(f\) 送到 \(m\mapsto(s\mapsto f(sm))\)。导出 \(F\) 得 \(\operatorname{Ext}_R^q(S,A)\)，其作用来自同一消解上的右乘预合成。Grothendieck 定理遂给出所写公式，模范畴均有足够内射对象。
\end{solution}

\begin{exercise}\label{b4:ex:15-change-rings-prime}
在上题中取 \(R=\mathbb Z\)、\(S=M=\mathbb F_p\)、\(A=\mathbb Z\)，其中 \(p\) 为素数。计算第二页及目标。
\end{exercise}
\begin{solution}
\(\operatorname{Hom}_{\mathbb Z}(\mathbb F_p,\mathbb Z)=0\)，而长度一自由消解给出 \(\operatorname{Ext}_{\mathbb Z}^1(\mathbb F_p,\mathbb Z)=\mathbb F_p\)，更高次为零。中间范畴是向量空间范畴，\(\operatorname{Ext}_{\mathbb F_p}^{p'}\) 在 \(p'>0\) 消失。因此第二页只有 \((0,1)\) 等于 \(\mathbb F_p\)，目标只在一次等于 \(\mathbb F_p\)。这里 \(p'\) 表示谱序列指标，以免与素数 \(p\) 混淆。
\end{solution}

\begin{exercise}\label{b4:ex:15-cyclic-coprime-kernel}
取扩张 \(1\to C_3\to C_6\to C_2\to1\)，系数 \(\mathbb F_2\) 作用平凡。利用群扩张谱序列证明膨胀 \(H^n(C_2,\mathbb F_2)\to H^n(C_6,\mathbb F_2)\) 对所有 \(n\ge0\) 为同构。
\end{exercise}
\begin{solution}
循环群公式给出正奇次 \(H^q(C_3,\mathbb F_2)=\mathbb F_2[3]=0\)，正偶次为 \(\mathbb F_2/3\mathbb F_2=0\)。所以第二页只在 \(q=0\) 留下 \(H^p(C_2,\mathbb F_2)\)。单行退化及边缘同态识别说明膨胀在各次为同构；次数零也是平凡作用下的恒等映射。
\end{solution}

\begin{exercise}\label{b4:ex:15-prime-square-transgression}
设 \(p\) 为素数，取 \(1\to C_p\to C_{p^2}\to C_p\to1\)，系数 \(\mathbb F_p\) 平凡。证明五项列中的传递映射为同构、最后的膨胀为零，不要求选择生成元后确定传递映射的符号。
\end{exercise}
\begin{solution}
五个候选项均为一维 \(\mathbb F_p\)。第一次膨胀把商群同态拉回，仍给出所有 \(C_{p^2}\to\mathbb F_p\) 的同态，所以为同构。限制到由 \(t^p\) 生成的核群时，任何同态都取值 \(pf(t)=0\)，所以限制为零。正合性使传递映射单；两端一维，故它为同构。最后一次膨胀的核是整个定义域，因此为零。
\end{solution}

\begin{exercise}\label{b4:ex:15-split-transgression}
设 \(G=N\rtimes Q\)，\(A\) 为平凡 \(G\)-模。证明五项列的传递映射 \(H^1(N,A)^Q\to H^2(Q,A)\) 为零。
\end{exercise}
\begin{solution}
每个一次类由同态 \(\phi:N\to A\) 表示，\(Q\)-不变表示 \(\phi(qnq^{-1})=\phi(n)\)。定义 \(f(n,q)=\phi(n)\)。半直积乘法给出
\(f((n,q)(n',q'))=\phi(n)+\phi(qn'q^{-1})=\phi(n)+\phi(n')\)，所以 \(f:G\to A\) 是同态，其限制为 \(\phi\)。限制因而满射，正合性使传递映射为零。这里使用了平凡系数假设，不能直接把同态公式用于任意交叉同态。
\end{solution}

\begin{exercise}\label{b4:ex:15-vanishing-range}
在 Grothendieck 定理的假设下，固定 \(M\)，并设 \(R^qF(M)=0\) 对 \(1\le q\le t\) 成立。证明边缘同态 \(R^nG(FM)\to R^n(GF)(M)\) 在 \(0\le n\le t\) 为同构。
\end{exercise}
\begin{solution}
在总次数 \(n\le t\) 上，第二页除底行外均为零。流入 \((n,0)\) 的第 \(r\ge2\) 页源为 \((n-r,r-1)\)：若横指标非负，则 \(1\le r-1<n\le t\)，对应已消失的行；若横指标负，也为零。没有出箭头，因此底行稳定。目标该次数只有一个非零相邻滤过商，故边缘同态为同构。
\end{solution}

\begin{exercise}\label{b4:ex:15-missing-acyclic-hypothesis}
设 \(p\) 为素数，令 \(F:\mathbb F_p\text{-}\mathbf{Mod}\to\mathbf{Ab}\) 为忘却函子，\(G=\operatorname{Hom}_{\mathbb Z}(\mathbb F_p,-)\)。说明 \(F,G\) 都左正合，但若删去内射像条件，所期待的复合函子谱序列会给出矛盾。
\end{exercise}
\begin{solution}
\(F\) 正合，且 \(GF(V)=V\) 作为交换群，因此复合函子正合，正次导出为零。若错误地套用谱序列，在 \(M=\mathbb F_p\) 处只可能有 \(q=0\) 行，而该行的 \(E_2^{1,0}=\operatorname{Ext}_{\mathbb Z}^1(\mathbb F_p,\mathbb F_p)=\mathbb F_p\) 非零。单行谱序列不可能再把它消去，与目标 \(R^1(GF)(M)=0\) 矛盾。失效条件恰好是：\(\mathbb F_p\) 在源范畴内射，忘却后却不对 \(G\) 零调。
\end{solution}
